\documentclass[11pt]{article}
\usepackage[margin=1.2in]{geometry}
\usepackage{amsmath,amssymb,amsthm}
\usepackage{hyperref}
\hypersetup{colorlinks=true, linkcolor=blue, urlcolor=blue}
\usepackage{parskip}
\emergencystretch=3em

\newcommand{\R}{\mathbb{R}}
\newcommand{\code}[1]{\texttt{#1}}

\title{Tide retrospective: singular composition, part 2\\
\large the point and locus theorems}
\author{Claude (Anthropic), with GPT-5.6 Sol consults}
\date{2026-08-10}

\begin{document}
\maketitle

\begin{abstract}
The composition arc opened by the adversarial audit of the note's
unnormalized degenerate theorem reaches its structural summit: the
point theorem \code{pencil\_families\_force\_germ\_eq\_at} --- if two
continuous nonnegative potentials, analytic at a common zero $p$,
produce Laplace families that agree beyond all orders against every
continuous compactly supported observable, then their germs at $p$
coincide --- and the locus theorem
\code{pencil\_families\_force\_eq\_near}, which assembles the
per-point conclusions into equality on an open neighborhood of the
whole zero set $W_0$. Two structural surprises. First, the assembly
needs no compactness: the union of the per-point neighborhoods is
already the required open set, so the audit's ``compactness point
selection'' step belongs to the note's contrapositive phrasing, not
to the mathematics, and what was planned as a third tide collapsed
into thirty lines of this one. Second, a genuine gap surfaced in the
other direction: the note quantifies its decay premise over
$\phi \in C_c^\infty$, while the Lean hypothesis takes all continuous
compactly supported $\phi$ --- a strictly stronger premise, since the
constructed observable uses the min/max bump, which is not smooth.
The theorems here are correct and are the engine; the smooth-class
upgrade is the one remaining step before the note's statement is
subsumed verbatim.
\end{abstract}

\section{Setting}

Part 1 (PR \#120) manufactured the endpoint's hypotheses: the
quadratic upper bound at a nonnegative $C^2$ zero, the least nonzero
diagonal of a nonvanishing analytic function with a normalized
direction witness, and a continuous bump. This tide composes them
with \code{analytic\_pencil\_difference\_not\_superpolynomial} into
statements a reader of the note can recognize.

\section{The thing formalised}

\begin{sloppypar}
The point theorem takes $L_1, L_2$ continuous, nonnegative, analytic
at $p$, vanishing at $p$, and the family premise: for every
continuous compactly supported $\phi$ and every $N$,
$t \mapsto \int \phi\,(e^{-tL_1} - e^{-tL_2})$ is $o(t^{-N})$. It
concludes $L_1 = L_2$ eventually in $\mathcal{N}(p)$. The proof runs
the audit's own reduction list: translate by $p$ (analyticity of the
shift composes; the power series of $L_2(p+\cdot) - L_1(p+\cdot)$ at
$0$ has a ball); by contradiction a witness where the germs differ
lands inside that ball (frequently-nonzero against a ball
neighborhood, with the $\varepsilon$-radius extracted from the
extended-real ball radius); the part-1 lemmas produce the least
diagonal $m$, the direction $x_0$ with $\|x_0\| = 3/2$, the quadratic
bound on $L_1(p+\cdot) + L_2(p+\cdot)$, and the bump; the endpoint
contradiction fires for the shifted pencil; and the Lebesgue
substitution \code{integral\_add\_left\_eq\_self} identifies the
endpoint's integral with the family integral at the observable
$(L_2 - L_1)\cdot\psi(\cdot - p)$. The locus theorem is then
\code{eventually\_nhds\_iff} + \code{choose} + a bi-union: no
compactness, no point selection.
\end{sloppypar}

\section{Where progress stalled}

Nowhere seriously --- the file passed its first full check after
three small repairs, all instances of catalogued classes.
\code{AnalyticAt.comp} wants the outer analyticity restated at the
shift's value at zero, and the \code{\(\triangleright\)} transport
produces a doubly-shifted point (write a \code{simpa}-typed
\code{have} instead). \code{integral\_add\_left\_eq\_self} takes the
function first and the translation second, mirroring its
multiplicative original; passing the point first silently elaborates
with the group set to the index type and dies asking for
\code{MeasurableSpace} on it. And the per-point goals of
\code{IsLittleO.congr'} arrive as applied lambdas, so
\code{beta\_reduce} precedes the rewrite --- the catalogued gotcha,
fired again.

\section{The observable-class gap}

The honest finding of the tide. The note's theorem premise is
``for every $\phi \in C_c^\infty$''; the Lean premise is ``for every
continuous compactly supported $\phi$''. A stronger premise makes a
weaker theorem: a reader holding the note's hypotheses cannot
instantiate the Lean one, because the observable built here is only
continuous. The repair is scoped in the tide log: extract an
analyticity radius at $p$, switch to a Mathlib \code{ContDiffBump}
whose closed support sits inside it, and glue
$(L_2 - L_1)\cdot\psi(\cdot - p)$ to a globally smooth function
pointwise (product of smooth functions on the analyticity ball,
locally zero off the support). The endpoint needs only continuity, so
only the observable-membership step changes. Until that lands, the
note's footnote must keep saying the reduction is partially outside
Lean.

\section{Mathlib lookup versus new mathematics}

\begin{sloppypar}
Lookup: \code{AnalyticAt.comp}, \code{AnalyticAt.contDiffAt},
\code{integral\_add\_left\_eq\_self},
\code{HasCompactSupport.comp\_homeomorph} with
\code{Homeomorph.subRight}, \code{eventually\_nhds\_iff},
\code{ENNReal.ofReal\_lt\_ofReal\_iff}. New mathematics: none --- and
the dissolution of the compactness step is a simplification the note
itself could adopt.
\end{sloppypar}

\section{What the next tide inherits}

The smooth observable-class upgrade
(\code{pencil\_families\_force\_germ\_eq\_at\_smooth} and its locus
form), whose premise matches the note verbatim; then the TeX
footnote rewrite, the \code{thm:singular} markers for the composed
theorems, and the pin bump.
\end{document}
